\section{Frozen source review and the actual comparison frontier}
\label{sec:review}

\subsection{Revision and review depth}

\begin{samepage}
The ProveIt review is frozen at
\begin{center}\small\ttfamily
8e76d515063c4b9fbd8743a846b457e3a01bc039.
\end{center}
\end{samepage}
The relevant recent commits are summarized below. Times in the source
metadata are UTC; these changes occurred during the evening of 26 October
in Pacific time. Source hashes and exact paths are recorded in the archive.

\begin{center}\small
\begin{tabular}{@{}llp{8.3cm}@{}}
\toprule Commit prefix & UTC time, 27 Oct. & Relevant change\\
\midrule
\texttt{6d3adc4506}&02:03&Eight round-2 packets placed.\\
\texttt{450252e9d8}&02:36&Five round-3 packets placed; omitted scripts restored.\\
\texttt{a7ae5ab55f}&03:15&Nine round-4 packets placed.\\
\texttt{7de9d09fec}&03:20&Public Anthropic Lean development linked.\\
\texttt{26bd83d4b3}&03:51&Five new affine/pairwise ZIP archives arrived.\\
\bottomrule
\end{tabular}
\end{center}

The project index contains 29 refinement packets: seven in round 1, eight
in round 2, five in round 3, and nine in round 4. The five newest incoming
archives raise the reviewed inventory to 34. We screened all indexed
claims and read the full mathematical sections that determine the current
frontier, inherited interfaces, and proposed order improvements. This is
not an assertion of an independent end-to-end proof audit of every earlier
packet. The inventory states which source claims are being summarized;
the main theorem has a narrower explicit dependency list.

During completion, commit \texttt{324f3d9b009a} placed the five incoming
packets into the project. All 213 placed files matched the reviewed ZIP
members by Git blob identity, including all 160 LaTeX, Python and JSON
files; the omitted files were documented ancillary or build files.
The closing observed main revision was \texttt{0cf10036e6ad}.
This placement changes the index to 34 packets without changing the
reviewed mathematics. The full identity check accompanies the archive.

\subsection{The five archives missing from the index}

All entries below concern the same generic normalized input model and
multiply $s/\sqrt\lambda$ in a base-two exponent. They are source claims
whose finite frontier certificates were inspected; the two leading
pairwise checkers were also replayed.

\begin{center}\small
\begin{tabular}{@{}p{6.1cm}rr@{}}
\toprule Incoming archive stem & Time coefficient & Space coefficient\\
\midrule
\texttt{sat\_pairwise\_certificates}&29.811609&29.811609\\
\texttt{sat\_affine\_contraction}&30.257036&25.989068\\
\texttt{sat\_affine\_geometry}&30.269873&29.487537\\
\texttt{sat\_affine\_cauchy}&30.413349&29.466030\\
\texttt{sat\_affine\_plane}&30.691338&29.466030\\
\bottomrule
\end{tabular}
\end{center}

The first two rows form the relevant time--space tradeoff among these
five. A smaller time coefficient does not erase the second row's smaller
memory coefficient. The best generic polynomial-space coefficient remains
$46.184845586731\ldots$ from \texttt{sat\_adaptive\_bounds}. The newest
pairwise packet separately gives $16.812233801883\ldots$ in polynomial
space at root degree at most $\ceil{s^{1/12}}$. That lower-degree result
has stronger assumptions and is not a competitor in the generic
degree-$s^{9/10}$ table.

The older indexed generic cached record was $30.891018920366\ldots$.
Comparing a new result only with that value would miss the incoming
$29.8116\ldots$ frontier. None of the 34 packets establishes a smaller
general order than~\eqref{eq:satorders}. The conditional improvements
require an additional algorithm or representation theorem that is not
provided by the corresponding packet.

\subsection{The two primary algorithms and formal verification}

The vendored OpenAI manuscript at upstream revision
\nolinkurl{7ef5c0cc5ca4b0e072b7a050fca09ab24bc6fa3a} proves the bounds
in~\eqref{eq:satorders}, with polynomial space, in its quantitative
appendix~\cite{openai}. Its published Lean scope page reports the weaker
general exponent $n(\log\log n)^2/\sqrt{\log n}$ for that formal
development~\cite{openaiscope}. The paper's sharper quantitative appendix
and its formalized theorem should not be conflated.

The Anthropic source is inspected at the user's exact commit
\begin{center}\small\ttfamily
1a6eee293b9e00cb432eef6018a21f523623a9af.
\end{center}
Its \texttt{MainResults.lean} states a finite-alphabet deterministic
Mathlib Turing-machine counting theorem, and \texttt{Solution.lean}
supplies proofs without extra hypotheses. Specializing to 3-SAT gives
\begin{equation}\label{eq:anthropic}
 2^{O\left(\min\left\{
       \frac{m}{\sqrt{\log m}},
       \frac{n\log\log n}{\sqrt{\log n}}
       \right\}\right)}.
\end{equation}
The stated counting form assumes normalized clauses and occurring named
variables; small cases and unused variables are handled by the usual
front end. Raw input normalization adds a polynomial factor in input
length. The separate ETH wrapper is also proved~\cite{anthropiccode}.

We retrieved the exact-commit upstream verification log. The
\texttt{verify (eth)} job reports successful Lean-kernel, nanoda, and
Comparator acceptance, with the deliberately unproved challenge statements
kept separate from \texttt{Solution.lean}~\cite{anthropicci}. This is
observed upstream verification, not a local rebuild performed here. The
older round-2 statement that the public development was unavailable is
therefore obsolete. The Anthropic development does not formalize the
ProveIt refinement packets or the new results in this report, and its
displayed time theorem does not supply the OpenAI paper's polynomial-space
claim.

\subsection{A stronger available matrix backend}

The newest pairwise packet still uses $\omega<19/8$ and dual exponent
$\alpha>3/10$. Current primary sources contain stronger results.
The OpenAI scope page records $\omega(\C)\le9/4$ and
$\alpha>0.465$~\cite{matrixscope}. Uzna\'nski's version-2 theorem gives
$\omega(1,1/2,1)=2$ and $\omega\le9/4$ over characteristic zero, with
positive exponent slack understood~\cite{uznanski}. We use these two
fixed shapes. This avoids assuming uniform constants for a continuously
varying family of rectangular shapes. The transfer to the exact modular
calculation is proved in Section~\ref{sec:matrix}.
